\documentclass[a4paper, 12pt]{article}

\usepackage[T1]{fontenc}
\usepackage[polish]{babel} 
\usepackage[utf8]{inputenc} 
\let\lll\undefined
\usepackage{setspace}
\usepackage{fancyhdr}
\usepackage{hyperref}
\usepackage{pdfpages}
\usepackage{listings}
\usepackage{color}
\usepackage{graphicx}
\usepackage{enumitem}
\usepackage{latexsym}
\pagestyle{fancy} 
\hypersetup{
    colorlinks=true,
    linkcolor=blue,
    filecolor=magenta,      
    urlcolor=cyan,
}
\usepackage{geometry}
\newgeometry{tmargin=2.5cm, bmargin=2.5cm, lmargin=2.5cm, rmargin=2.5cm} 
\newcommand{\mainmatter}{\clearpage \cfoot{\thepage\ of \pageref{LastPage}}
\pagenumbering{arabic}} 
\lstset{language=bash}  
\begin{document}

	\begin{titlepage}
\includegraphics[width = 40mm]{logo.jpg}
		\begin{center}
    			\vspace{3cm}
    					\Large\textit{\textbf{Procesy biznesowe (definicje, kryteria i klasyfikacje, przykłady)}}
   			\vspace{4cm}
		\end{center} 

		\hfill\begin{minipage}{0.54\textwidth}
			\Large Wykonali:\newline
				1. Ivan Prakapets 295139 \newline
				2. Aliaksandr Karolik 295138
		\vspace{\baselineskip}
		\end{minipage}
		
		\hfill\begin{minipage}{0.54\textwidth}
			\Large Sprawdzający:\newline
		 		prof.dr hab.inż.Andrzej Dzieliński
\vspace{\baselineskip}
		\end{minipage}

		\hfill\begin{minipage}{0.54\textwidth}
			\Large Data:10.11.2019
			\vspace{\baselineskip}
		\end{minipage}
	\end{titlepage}
\newpage
\mainmatter
\setlength{\headheight}{15pt}
\doublespacing
\tableofcontents
\newpage

\linespread{0.5}
\setlist{nolistsep}

\section{Wstęp}
\hspace*{1.5 cm}Poniższy dokument jest  pracą semestralną na temat "Procesy biznesowe (definicje, kryteria i klasyfikacje, przykłady)". W tym dokumencie postaramy się zdefiniować co to jest proces biznesowy oraz zbadać jakie są kryteria i klasyfikacji procesów biznesowych. Również postaramy się rozpatrzyć kilka procesów biznesowych w różnych firmach aby na ich podstawie zaobserwować różne rodzaje procesów biznesowych. 
\section{Definicja procesu biznesowego}
\hspace*{1.5 cm}Poniżej zostaną zaprezentowane różne definicje pojęcia proces biznezosy:

\hspace*{1.5 cm}Proces biznesowy jest zbiorem czynności, ma jeden lub więcej rodzajów wejść i tworzy wartość wyjściową dla klienta. Proces biznesowy posiada swój cel, a oddziałują na niego zdarzenia zachodzące w świecie zewnętrznym lub w innych procesach.
(Hammer and Champy 1993)

\hspace*{1.5 cm} Proces biznesowy jest to zbiór powiązanych procedur lub działań, które wspólnie zapewniają osiągnięcie celu biznesowego lub celu polityki, zwykle w ramach struktury organizacyjnej definiującej funkcjonalność ról i zależności pomiędzy nimi. (Workflow Management Coalition 1999) 

\hspace*{1.5 cm}Proces biznesowy może być postrzegany jako struktura czynności zaprojektowanych na działania nakierowane na klienta końcowego i dynamiczne zarządzanie przepływami związanymi z produktami, informacją, środkami finansowymi, wiedzą i wizją.(Stock and Lambert 2001)        

\hspace*{1.5 cm}Proces biznesowy lub inaczej metoda biznesowa, która funkcjonuje w każdym przedsiębiorstwie. Są to zadania ze sobą powiązane, które prowadzą do osiągnięcia wyznaczonego efektu. Najważniejszym celem tego procesu jest zrozumieć klienta, dostawców i słabe strony tego procesu.(A. Bitkowska 2009, s.26)

\hspace*{1.5 cm}Określony zbiór czynności biznesowych, które stanowią niezbędne kroki w celu osiągnięcia celu biznesowego. Obejmuje on przepływ i użycie informacji oraz zasobów.(Object Management Group – BPMN v. 2.0 2011) 
\subsection{Podsumowanie czym jest proces biznesowy} 
\hspace{1.5 cm} Bazując na powyższych definicjach, można przyjąć, że proces biznesowy, to zbiór powiązanych ze sobą czynności, które przekształcają wejścia w wyjścia według określonych reguł, w oparciu o określone zasoby i w efekcie prowadzą do dostarczenia klientowi produktu/usługi realizując tym samym cele biznesowe organizacji.
\section{Kryteria i klasyfikacje procesów biznesowych}
\subsection{Rodzaje kryteriów}
\subsection{Rodzaje klasyfikacji}
\hspace*{1.5 cm}W litera-turze  przedmiotu  spotkać  można  wiele  różnych  kryteriów  podziału  i klasyfikacji  procesów.  I  tak,  według  M.  Portera  [Porter,  1985,  s.  23]  wy-różnić  można  dwa  podstawowe  rodzaje  procesów:  podstawowe  i  po-mocnicze. Do procesów podstawowych zaliczył on: 
\begin{itemize}
	\item Logistykę „na wejściu”, która obejmuje działania związane z przygo-towaniem produkcji,
	\item wytwarzanie produktu, 
 	\item logistykę  „na  wyjściu”,  która obejmuje  działania związane ze sprze-dażą,
	\item marketing,
	\item usługi posprzedażne.
\end{itemize}

\hspace*{1.5 cm}Natomiast wśród pomocniczych wymienił procesy związane z: 
\begin{itemize}
	\item Zarządzaniem całą jednostką,
	\item zarządzaniem zasobami ludzkimi,
	\item zaopatrzeniem,
	\item rozwojem mającym na celu doskonalenie produktów i procesów.
\end{itemize}

\hspace*{1.5 cm} Z  kolei  R.S.  Kaplan  i  R.  Cooper  [Kaplan,  Cooper,  2001, s.  99],  [Kaplan, Cooper, 2000, s. 200] wyróżnili procesy:
\begin{itemize}
	\item Innowacyjne – związane z określaniem rynku docelowego oraz two-rzeniem oferty produktu (usługi), 
	\item operacyjne  –  dotyczące  wytwarzania  produktu  (usługi)  i  dostarcza-nia go klientowi,
	\item obsługi posprzedażnej – obejmujące obsługę klienta po dostarczeniu mu produktu,
\end{itemize}
\hspace*{1.5 cm} Oraz podzielili je na:
\begin{itemize}
	\item Konieczne – których wykonanie jest niezbędne do dostarczenia war-tości  i  których  nie  można  obecnie  poprawić,  uprościć,  zredukować czy wyeliminować, 
	\item istotne – dostarczające wartość, aczkolwiek możliwe jest ich uprosz-czenie i poprawienie,
	\item nieistotne – które powinny być wyeliminowane. 
\end{itemize}
\hspace*{1.5 cm} Z kolei najobszerniejszą klasyfikację procesów przedstawiła organizacja APQC (American Productivity Quality Center), która przygotowała Model Klasyfikacji Procesów (z ang. Process Classification Framework PCF). Wspomniana  organizacja  zaproponowała  wprowadzenie  12  ka-tegorii procesów, które podzielone zostały na dwie grupy:
\begin{itemize}
	\item Procesy  operacyjne  –  traktowane  jako  kluczowe,  dla  danego  pod-miotu gospodarczego, 
	\item procesy  wspomagające  –  stanowiące  wsparcie  dla  procesów  opera-cyjnych.
\end{itemize}
\hspace*{1.5 cm} Do  pierwszej  grupy  zaliczono  5  kategorii  charakteryzujących  działalność  podmiotu  decydujących  o  zakresie  funkcjonowania przedsiębiorstwa,  natomiast  do  procesów  wspomagających  zakwalifikowano  7 kategorii,  które  przenikają  wszystkie  procesy  podstawowe  i  jednocześnie są podstawą do wydzielania ich na zewnątrz (outsourcing). Ilustrację  graficzną  klasyfikacji  procesów  zgodnej  z  modelem  APQC  przed-stawia tablica 1.

\begin{table}[h]
	\begin{center}
		\scalebox{1}{\begin{tabular}{ | l | l | }
				\hline
				Proces & Charakterystyka \\ \hline
				Operacyjne & 1.0 - Opracowanie wizji i strategii \\
						   & 2.0 - Rozwój i zarządzanie produktami i usługami  \\
						   & 3.0 -  Marketing i sprzedaż produktów i usług \\ 
						   & 4.0 - Zaopatrzenie, realizacja i dostawa produktów/usług \\
						   & 5.0 - Zarządzanie obsługą klienta		   \\ \hline
				Wspomagające & 6.0 -Organizacja i zarządzanie 	kapitałem ludzkim \\
				             & 7.0 - Zarządzanie technologią informatyczną\\
				             & 8.0 - Zarządzanie zasobami finansowymi\\
				             & 9.0 - Nabywanie, budowa i zarządzanie mieniem \\
				             & 10.0 - Zarządzanie ochroną środowiska oraz bezpieczeń-stwem i higieną pracy \\
				             & 11.0 - Zarządzanie relacjami zewnętrznymi\\
				             & 12.0 -  Zarządzanie wiedzą, doskonaleniem i zmianą\\
				\hline
		\end{tabular}}
	\end{center}
	\caption{\label{tab:bolts} Tablica 1. Klasyfikacja procesów według modelu APQC}
\end{table}
\hspace*{1.5 cm}Model APQC nie ogranicza się do przedstawienia listy 12 kategorii procesów,  rozwija  się  proponując  następującą  zależność:  kategoria  procesów  –  grupa  procesów  –  procesy  –  czynności  (działania).
\subsection{Podsumowanie}

\section{Przykłady biznes procesów w różnych rodzajach firm}
\subsection{Prykład biznes procesu w firmie IT}
\subsubsection{Opis ogólny procesu}
\hspace*{1cm} Potrzebujemy zakupić nowe serwery z potrzeby zaspokojenia bieżącego zapotrzebowania na moc obliczeniową, przyrostu danych oraz wymiana obecnie posiadanego systemu na nowszy, bo jest brak kompatybilności z nowymi systemami operacyjnymi, brak wsparcia producenta, z tego wynika wysokie ryzyko awarii. Poniżej jest dokładne opisane problemy, informacje, wykresy oraz charakterystyki dotyczące obecnych i nowych serwerów.\newline
\subsubsection{Występujące problemy w firmie} 
		\begin{itemize}
		        \item \texttt{Na obecnym serwerze baza danych wzrasta się około 100 Gb co rok. Ze względu na duże przyrosty danych na obecnych serwerach brakuje miejsca. Brak przestrzeni na dane powoduje coraz wolniejsze działanie serwera. W przypadku awarii niemożliwym staje się szybka naprawa. W obecnej postój pracy wynosi 5 dni.}
	        	\item \texttt{Problemy z komunikacją oraz transferem danych między serwerami spowodowana przestarzałością sprzętu}
	        	\item \texttt{Niski poziom Uptime, poziom dostępności systemu określony w procentach. Informujący o czasie ciągłego i bezawaryjnego działania serwerów w ciągu roku. Można spodziewać się awarii w niedługim okresie.}
	        	\item \texttt{Okres gwarancji na sprzęt oraz oprogramowanie wygasł, wiąże się to z częstrzą możliwością występowania corac cięższych awarii.}
	        	\item \texttt{Przestarzały sprzęt wycofany już z produkcji i powszechnego użytku. Posiadany sprzęt zurzywa się i traci pierwotą niezawodność. Generuje to ryzyko braku możliwości wymiany części i podzespołów w przypadku uszkodzeń i awarii.}
	        	\item \texttt{ Stare oprogramowanie traci swoją kompatybilność z nowymi systemami i ma słabsze możliwości względem nowych.}
	        	\item \texttt{Niewystarczająca moc obliczeniowa powoduje powolną pracę serwerów. Dane przesyłane do aplikacji są coraz wolniej selekcjonowane oraz generowanie z długim czasem oczekiwania na wyniki. Wpływ na to ma szybkość procesowa, pamięć RAM, pojemność dysku, przepustowość, itp.}
		    \end{itemize}
	  \subsubsection{Wnioski} 
	
			 	\hspace*{1cm} Pozostaje mało miejsca i istnieje potrzeba rozszerzenia przestrzeni roboczej poprzez podłączenie dodatkowych urządzeń pamięci masowej lub zwiększenie pojemności istniejącej macierzy dyskowej. Niestety uniemożliwia to przestarzały sprzęt niekompatybilny z odpowiednikami dostępnymi na rynku.\newline
			\hspace*{1cm}  Wszystkie powyższe problemu powodują przerwy, awarie oraz  przestoje - \textbf{generuje to straty materialne oraz oszczerbki na reputacji firmy.}
			
			\newpage
			\subsubsection{Rozwiązanie problemów}
				TUTAJ COS JESZCZE
				\subsubsection{Klastry SQL Server 2019} 	
	\hspace*{1cm} Klastry to technologia, która pozwala łączyć kilka systemów komputerowych w całość. 
Cel klastra to z duplikowanie  serwerów i zapewnienie nieprzerwanego funkcjonowania aplikacji. W serwerach bazodanowych wszystkie informacje są zapamiętywane na każdej równolegle pracującej maszynie. W razie awarii systemy i urządzenia łączą się działającym serwerem, zawierającym te same dane.
\newline
	\hspace*{1cm}Klastry SQL Server 2019 to nowy sposób wykorzystania SQL Server do tworzenia wartościowych baz relacyjnych i przechowywania dużych zbiorów danych na jednej, skalowalnej platformie. Dzięki temu analiza danych oraz aplikacje łączące się z bazą są bardziej elastyczne i działają wydajniej. Duże klastry danych SQL Server 2019 zapewniają kompletną platformę i pomagają zwiększyć sukces  organizacji.
\subsubsection{Opis przepływu procesu biznesowego}

\subsection{Przykład biznes procesu w firmie marketingowej}
\subsubsection{Ogólny opis procesu}
\subsubsection{Problemy które proces rozwiązuje}
\subsubsection{Opis przepływu procesu biznesowego}

\section{Wnioski} 


\label{LastPage}~
\label{LastPageOfBackMatter}~		
\end{document}